\documentclass{article}

% ---------- packages ----------
\usepackage{amsmath}
\usepackage{amsthm}
\usepackage{xcolor}
\usepackage{etoolbox}
\usepackage{hyperref}
\usepackage{cleveref}
\usepackage{xparse}
\usepackage[most]{tcolorbox}

% ---------- theorem styles ----------
\newtheoremstyle{mytheoremstyle}{}{}{}{}{\sffamily\bfseries}{.}{ }{}
\newtheoremstyle{myconjecturestyle}{}{}{}{}{\sffamily\bfseries}{.}{ }{\thmnote{#3}}

% ---------- theorem environments ----------
\theoremstyle{mytheoremstyle}{\newtheorem{definition}{Definition}[section]}
\theoremstyle{mytheoremstyle}{\newtheorem{proposition}[definition]{Proposition}}
\theoremstyle{mytheoremstyle}{\newtheorem{theorem}[definition]{Theorem}}
\theoremstyle{mytheoremstyle}{\newtheorem{lemma}[definition]{Lemma}}
\theoremstyle{mytheoremstyle}{\newtheorem{corollary}[definition]{Corollary}}

\tcbset{
    tbox_DPLTC_style/.style={
        enhanced jigsaw,
        colback=#1!10, colframe=#1!80!black,
        boxrule=0pt,
        fonttitle=\sffamily\bfseries, coltitle=black,
        separator sign={}, label separator={},
        description font=\normalfont\sffamily,
        description delimiters={(}{)},
        attach title to upper, after title={.\ },
        frame hidden, borderline west={2pt}{0pt}{#1},
        sharp corners,
        top=2pt, bottom=2pt, left=5pt, right=5pt,
        beforeafter skip=10pt, breakable,
    },
    tbox_DPLTC_style*/.style={
        tbox_DPLTC_style=#1,
        interior hidden,
        top=-2pt, bottom=-2pt,
    },
}

% ---------- colours ----------
\definecolor{tcol_DEF}{HTML}{E40125}
\definecolor{tcol_PRP}{HTML}{EB8407}
\definecolor{tcol_LEM}{HTML}{05C4D9}
\definecolor{tcol_THM}{HTML}{1346E4}
\definecolor{tcol_COR}{HTML}{7904C2}

\tcolorboxenvironment{definition}{tbox_DPLTC_style=tcol_DEF}
\tcolorboxenvironment{proposition}{tbox_DPLTC_style=tcol_PRP}
\tcolorboxenvironment{theorem}{tbox_DPLTC_style=tcol_THM}
\tcolorboxenvironment{lemma}{tbox_DPLTC_style=tcol_LEM}
\tcolorboxenvironment{corollary}{tbox_DPLTC_style=tcol_COR}
\tcolorboxenvironment{proof}{boxrule=0pt,boxsep=0pt,blanker,borderline west={2pt}{0pt}{CadetBlue!80!white},left=8pt,right=8pt,sharp corners,before skip=10pt,after skip=10pt,breakable}
\tcolorboxenvironment{remark}{tbox_DPLTC_style*=tcol_REM}
\tcolorboxenvironment{remarks}{tbox_DPLTC_style*=tcol_REM}
\tcolorboxenvironment{example}{tbox_DPLTC_style*=tcol_EXA}
\tcolorboxenvironment{examples}{tbox_DPLTC_style*=tcol_EXA}
\tcolorboxenvironment{cthm}{boxrule=0pt,boxsep=0pt,colback={gray!10},left=8pt,right=8pt,enhanced jigsaw, borderline west={2pt}{0pt}{gray},sharp corners,before skip=10pt,after skip=10pt,breakable}

% ---------- cleveref hooks (TeXLive 2025+) ----------
\AddToHook{env/proposition/begin}{\crefalias{definition}{proposition}}
\AddToHook{env/lemma/begin}{\crefalias{definition}{lemma}}
\AddToHook{env/theorem/begin}{\crefalias{definition}{theorem}}
\AddToHook{env/corollary/begin}{\crefalias{definition}{corollary}}

\crefformat{definition}{\normalfont\sffamily\bfseries\color{tcol_DEF}#2definition~#1#3}
\Crefformat{definition}{\normalfont\sffamily\bfseries\color{tcol_DEF}#2Definition~#1#3}
\crefformat{proposition}{\normalfont\sffamily\bfseries\color{tcol_PRP}#2proposition~#1#3}
\Crefformat{proposition}{\normalfont\sffamily\bfseries\color{tcol_PRP}#2Proposition~#1#3}
\crefformat{lemma}{\normalfont\sffamily\bfseries\color{tcol_LEM}#2lemma~#1#3}
\Crefformat{lemma}{\normalfont\sffamily\bfseries\color{tcol_LEM}#2Lemma~#1#3}
\crefformat{theorem}{\normalfont\sffamily\bfseries\color{tcol_THM}#2theorem~#1#3}
\Crefformat{theorem}{\normalfont\sffamily\bfseries\color{tcol_THM}#2Theorem~#1#3}
\crefformat{corollary}{\normalfont\sffamily\bfseries\color{tcol_COR}#2corollary~#1#3}
\Crefformat{corollary}{\normalfont\sffamily\bfseries\color{tcol_COR}#2Corollary~#1#3}

% ---------- \Nref implementation ----------
\makeatletter

\let\orig@thmhead\thmhead
\renewcommand{\thmhead}[3]{%
  \gdef\@currentthmname{#3}%
  \gdef\@currentthmtype{#1}%
  \orig@thmhead{#1}{#2}{#3}%
}

\NewDocumentCommand{\namedlabel}{m}{%
  \ifx\@currentthmname\@empty
    \PackageError{Nref}{%
      \string\namedlabel\space used on label '#1', but this theorem has no name%
    }{%
      Add an optional name: \string\begin{theorem}[Name]\string\namedlabel{#1}%
    }%
  \else
    \protected@write\@auxout{}{%
      \string\newlabel{thmname@#1}{{\@currentthmname}}%
    }%
    \label{#1}%
  \fi
}

\newcommand{\Nref@format}[2]{%
  % #1 = env type, #2 = name
  \normalfont\sffamily\bfseries
  \ifstrequal{#1}{Definition}{\color{tcol_DEF}}{%
  \ifstrequal{#1}{Proposition}{\color{tcol_PRP}}{%
  \ifstrequal{#1}{Lemma}{\color{tcol_LEM}}{%
  \ifstrequal{#1}{Theorem}{\color{tcol_THM}}{%
  \ifstrequal{#1}{Corollary}{\color{tcol_COR}}{%
  }}}}}%
  #2%
}

% Helper: receives the \newlabel value as two brace groups and formats them
\newcommand{\Nref@unpack}[2]{\Nref@format{#2}{#1}}

\NewDocumentCommand{\Nref}{m}{%
  \ifcsname r@thmname@#1\endcsname
    \hyperref[#1]{%
      \expandafter\Nref@unpack\csname r@thmname@#1\endcsname
    }%
  \else
    \PackageError{Nref}{%
      \string\Nref\space used on label '#1' but no name was stored%
    }{%
      Use \string\namedlabel{#1} instead of \string\label{#1} inside the theorem%
    }%
  \fi
}

\makeatother

% ---------- document ----------
\begin{document}

\section{Definitions and theorems}

\begin{definition}[Topological Space]
  \namedlabel{def:topspace}
  A \emph{topological space} is a set $X$ together with a collection $\tau$
  of subsets of $X$ satisfying the usual axioms.
\end{definition}

\begin{definition}
  \label{def:continuous}
  A map $f \colon X \to Y$ between topological spaces is \emph{continuous}
  if the preimage of every open set is open.
\end{definition}

\begin{proposition}[Composition of Continuous Maps]
  \namedlabel{prop:composition}
  If $f \colon X \to Y$ and $g \colon Y \to Z$ are continuous, then
  $g \circ f \colon X \to Z$ is continuous.
\end{proposition}

\begin{proposition}
  \label{prop:unnamed}
  The identity map on any topological space is continuous.
\end{proposition}

\begin{lemma}[Tube Lemma]
  \namedlabel{lem:tube}
  Let $X$ be a topological space and $K$ compact. If $W$ is open in
  $X \times K$ and contains $\{x_0\} \times K$, then $W$ contains a tube
  $U \times K$ for some open $U \ni x_0$.
\end{lemma}

\begin{lemma}
  \label{lem:unnamed}
  A finite union of compact sets is compact.
\end{lemma}

\begin{theorem}[Tychonoff's Theorem]
  \namedlabel{thm:tychonoff}
  An arbitrary product of compact spaces is compact in the product topology.
\end{theorem}

\begin{theorem}
  \label{thm:unnamed}
  A continuous image of a compact space is compact.
\end{theorem}

\begin{corollary}[Extreme Value Theorem]
  \namedlabel{cor:evt}
  A continuous real-valued function on a compact space attains its maximum
  and minimum.
\end{corollary}

\begin{corollary}
  \label{cor:unnamed}
  The continuous image of a connected space is connected.
\end{corollary}

\section{References}

Below we exercise every combination of named/unnamed and \verb|\Nref|/\verb|\Cref|.

\subsection*{Named environments}

\begin{itemize}
  \item \Nref{def:topspace} (definition, named, \texttt{\string\Nref})
  \item \Cref{def:topspace} (definition, named, \texttt{\string\Cref})
  \item \Nref{prop:composition} (proposition, named, \texttt{\string\Nref})
  \item \Cref{prop:composition} (proposition, named, \texttt{\string\Cref})
  \item \Nref{lem:tube} (lemma, named, \texttt{\string\Nref})
  \item \Cref{lem:tube} (lemma, named, \texttt{\string\Cref})
  \item \Nref{thm:tychonoff} (theorem, named, \texttt{\string\Nref})
  \item \Cref{thm:tychonoff} (theorem, named, \texttt{\string\Cref})
  \item \Nref{cor:evt} (corollary, named, \texttt{\string\Nref})
  \item \Cref{cor:evt} (corollary, named, \texttt{\string\Cref})
\end{itemize}

\subsection*{Unnamed environments (\texttt{\string\Cref} only)}

\begin{itemize}
  \item \Cref{def:continuous}
  \item \Cref{prop:unnamed}
  \item \Cref{lem:unnamed}
  \item \Cref{thm:unnamed}
  \item \Cref{cor:unnamed}
\end{itemize}

\subsection*{Inline prose}

By \Nref{thm:tychonoff}, products of compact spaces are compact.
Combined with \Nref{lem:tube}, this yields \Cref{cor:evt}.
The unnamed \Cref{thm:unnamed} gives compactness of images,
and \Cref{def:continuous} provides the relevant notion of continuity.

% \Nref{def:continuous}

\end{document}